\documentclass[12pt]{article}
\usepackage[utf8]{inputenc}
\usepackage[T1]{fontenc}
\usepackage{hyperref}
\title{Uru-Eu-Wau-Wau Language Resources for K-12 Curriculum Development}
\author{Compiled by Claude Sonnet 5 (claude-sonnet-5) for Andr\'as Kornai}
\date{}
\begin{document}
\maketitle
\section{Introduction}
Uru-Eu-Wau-Wau (ISO 639-3 urz) is a Kagwahiva (Tupi-Guarani) variety spoken in Rond\^onia, Brazil. Speaker estimates in the sources found vary widely, from about 60 to about 170. Search turned up very little: no primer, textbook or dictionary specific to this variety. This report is therefore short.
\section{Primary Educational Resources}
None found in search results.
\section{Linguistic and Reference Materials}
\begin{itemize}
\item Wikipedia entry on the Uru-Eu-Wau-Wau dialect (general overview, phoneme inventory, not verified against primary sources): \url{https://en.wikipedia.org/wiki/Uru-Eu-Wau-Wau_dialect}
\item Wikipedia entry on the Uru-Eu-Wau-Wau people: \url{https://en.wikipedia.org/wiki/Uru-Eu-Wau-Wau}
\item Betts, LaVera (2012), reference catalogued by Glottolog under dictionaries and vocabularies; it appears to be based on brief 1991 contacts, so probably a small wordlist at most: \url{https://glottolog.org/resource/reference/id/309049}
\item Article in Revista Língua Viva (Portuguese), a general introduction to the people and language: \url{https://periodicos.unir.br/index.php/linguaviva/article/download/1133/1229/4102}
\end{itemize}
\section{Religious and Community Resources}
None found in search results.
\section{Digital and Audio Resources}
Endangered Languages Project entry (status and basic information only): \url{https://endangeredlanguages.com/node/86633}
\section{Curriculum Development Context}
The language is not known to be taught in schools. Any curriculum would have to begin with new documentation made together with the community. Related Kagwahiva documentation (for example an ELAR collection mentioned by Wikipedia) should be checked, but I did not see its URL in search results and could not confirm its contents.
\end{document}